\documentclass{article}

\usepackage{fontsize}
\changefontsize{14}

\usepackage[a4paper, top=1.5cm, bottom=1.5cm, left=2.5cm, right=1.5cm]{geometry}

\usepackage[T2A]{fontenc}   % кодировка шрифтов для кириллицы
\usepackage[utf8]{inputenc} % кодировка исходного файла (UTF-8)
\usepackage[english, russian]{babel} % правила переносов и типографики

\usepackage{romannum} % римские цифры

\usepackage{amsmath} % вставки текста в формулы, нумерация формул


\usepackage{siunitx} %значок градуса

\usepackage{makecell} %переносы в таблице
\usepackage{graphicx} %вставка изображений



\begin{document}

\pagenumbering{arabic}


%ТИТУЛЬНИК


\thispagestyle{empty}

\begingroup
\setlength{\parindent}{0pt}

\begin{bfseries}

\vspace*{0cm}
Министерство образования и науки РФ

\vspace{2cm}
Московский государственный технический университет имени\\Н.Э. Баумана

\vspace{2cm}
Факультет «Фундаментальные науки»

\vspace{2cm}
Кафедра «Теоретическая механика» имени профессора\\ Н.Е. Жуковского

\vspace{2cm}

Домашнее задание № 4\\
Статика

\end{bfseries}

\vspace{2cm}

\textbf{Вариант} \hspace{0.5cm} 20

\vspace{0.5cm}

\textbf{Группа} \hspace{0.5cm} СМ11-21Б

\vspace{0.5cm}

\textbf{Студент} \hspace{0.5cm} Усачев Е.М. \hspace{1.5cm} \underline{\hspace{3cm}} 
\hspace{1cm}\underline{\hspace{3cm}}\\
\small\hspace*{8.5cm} подпись \hspace{2.7 cm} дата \normalsize

\vspace{0cm}
\textbf{Преподаватель} \hspace{0.5cm} Белый В.С.


\vfill

\vspace{1cm}
\begin{center}
\textbf{Москва 2026 год}
\end{center}

\endgroup

%ТИТУЛЬНИК


\newpage

\begin{center}\large \textbf{Плоская статика 1} \end{center}

\noindent Дано:\\
$AB=BC=CD=l$; $\alpha = 45^\circ$; $\beta = 60^\circ$; $\lambda = 30^\circ$; $C_{упр} = Pl$; $M = Pl$;  $q_m = \frac{2P}{l}$.

\medskip

\noindent Определить реакции опор $A$ и $B$.

\bigskip 
\bigskip 

\begin{center}\includegraphics[width=0.8\textwidth]{scheme1.png}\\
\text{Рис. 1. Схема системы 1} \end{center}


\newpage



Данная система является статически неопределимой, поэтому разобьем ее на 3 системы (стержни $AB$, $BC$, $CD$), добавив силы, действующие в их соединениях (в шарнирах $B$ и $C$).

Рассмотрим стержень $AB$. Запишем для него уравнения равновесия:
\begin{equation}
\sum F_x = X_B + X_A - Psin60^\circ = 0;
\end{equation}
\begin{equation}
\sum F_y = Y_A - Y_B - Pcos60^\circ = 0;
\end{equation}
\begin{equation}
\sum M(B) = X_Al - Psin60^\circ\frac{l}{2} = 0.
\end{equation}

Из уравнения $(3)$ найдем горизонтальную проекцию реакции опоры $A$ $X_A$:
\begin{equation}
X_A = 0,5Psin60^\circ = 0,43P.
\end{equation}

Подставив уравнение $(4)$ в уравнение $(1)$ получим:
\begin{equation}
X_B = Psin60^\circ - 0,5Psin60^\circ = 0,5Psin60^\circ = 0,43P.
\end{equation}

Рассмотрим стержень $BC$. Запишем для него уравнения равновесия:
\begin{equation}
\sum F_x = X_C - X_B = 0;
\end{equation}
\begin{equation}
\sum F_y = Y_B + Y_L+Y_C - Q = 0;
\end{equation}
\begin{equation}
\sum M(C) = M_{ynp} - \frac{1}{3}Ql + 0,5Y_Ll + Y_Bl = 0.
\end{equation}

Из уравнения $(6)$:
\begin{equation}
X_C = X_B = 0,5Psin60^\circ.
\end{equation}

Пружина в шарнире $C$ заведена на угол $\lambda$ по часовой стрелке, поэтому она создает вращательный момент $M_{ynp}$, который стремится повернуть стержень $BC$ против часовой стрелки:
\begin{equation}
M_{ynp} = C_{ynp}\frac{\pi}{6} = \frac{\pi}{6}Pl.
\end{equation}

$Q$ - это распределенная сила, действующая на стержень $BC$ в точке на расстоянии $\frac{2}{3}l$ от точки $B$, которая равна:
\begin{equation}
Q = 0,5q_ml = 0,5l\frac{2P}{l} = P.
\end{equation}

Для стержня $CD$ запишем только уравнение моментов относительно точки $D$:
\begin{equation}
\sum M(D) = M - M_{ynp} + X_Clcos30^\circ - Y_Clsin30^\circ = 0.
\end{equation}

Из этого уравнения выразим $Y_C$ и подставим вместе с уравнением $(9)$ уравнение $(7)$:
\begin{equation}
Y_C = \frac{Pl - \frac{\pi}{6}Pl + X_Clcos30^\circ}{lsin30^\circ} = \frac{\frac{5\pi}{6}P + 0,5Psin60^\circ cos30^\circ}{sin 30^\circ} = 5,99P;
\end{equation}
\begin{equation}
Y_B + Y_L +5,99P -P = 0.
\end{equation}

Из уравнения $(8)$ выразим силу реакции шарнирного стержня в точке $L$ $Y_L$ и подставим его в уравнение $(14)$:
\begin{equation}
Y_L = \frac{\frac{1}{3}Pl - \frac{\pi}{6}Pl - Y_Bl}{0,5l} = -0,38P -2Y_B;
\end{equation}
\begin{equation}
Y_B - 0,38P - 2Y_B +5,99P -P = 0;
\end{equation}
\begin{equation}
Y_B = 4,61P.
\end{equation}

Подставим уравнение $(17)$ в уравнение $(2)$:
\begin{equation}
Y_A - 4,61P - Pcos60^\circ = 0;
\end{equation}
\begin{equation}
Y_A = 4,61P + 0,5P = 5,11P.
\end{equation}

Тогда силы реакции опор $A$ и $B$ равны:
\begin{equation}
R_A = \sqrt{X^2_A + Y^2_A} = \sqrt{(0,43P)^2 + (5,11P)^2} = 5,13P;
\end{equation}
\begin{equation}
R_B = \sqrt{X^2_B + Y^2_B} = \sqrt{(0,43P)^2 + (4,61P)^2} = 4,63P.
\end{equation}

Так как сила $Y_L$ получилась отрицательной в ходе расчетов, то она будет направлена в противоположную сторону. Остальные силы направлены верно.

\medskip
Вывод: в результате расчетов были получены значения сил реакции опор $A$ и $B$, которые представлены в таблице 1.

\medskip


\begin{flushright}
  \textit{Таблица 1} \hspace*{0.4cm}
\end{flushright}

\vspace*{-0.9cm}

\begin{center}
  Значения сил реакции опор
\end{center}

\begin{tabular}{|c|c|c|c|c|}
  \hline
  \rule{0pt}{1cm} № & \hspace*{0.5cm} Название величины \hspace*{0.5cm} & \makecell{Обозначение\\ величины} & Значение & \makecell{Единицы\\ измерения} \\[0.5cm]
  \hline
  \rule{0pt}{1cm} 1  & \makecell{Сила реакции опоры $A$} & $R_A$ & $5,13P$ & Н \\[0.5cm]
  
  \hline
  \rule{0pt}{1cm} 2  & \makecell{Сила реакции опоры $B$} & $R_B$ & $4,63P$ & Н \\[0.5cm]
  \hline
\end{tabular}

\newpage


\begin{center}\large \textbf{Пространственная статика 1} \end{center}

\noindent Дано:\\
$P$; $l$; $R = 2r = l$; $l_1 = 2l$; $AB = 2l$; $\varphi= \beta = 60^\circ$; $\alpha = 30^\circ$; $M = 0,5Pl$.

\medskip

\noindent Определить реакции подпятника $A$ и подшипника $B$, а также необходимую для равновесия силу $Q$.

\bigskip 
\bigskip 

\begin{center}\includegraphics[width=0.8\textwidth]{scheme3.png}\\
\text{Рис. 3. Схема системы 2} \end{center}


\newpage

Данная система является статически определимой. Запишем для нее условия равновесия:
\begin{equation}
\sum F_x = X_A -Pcos60^\circ + Qcos30^\circ sin60^\circ + X_B = 0;
\end{equation}
\begin{equation}
\sum F_y = Y_A - Psin60^\circ - Qsin30^\circ + Y_B = 0;
\end{equation}
\begin{equation}
\sum F_z = Z_A - Qcos30^\circ cos60^\circ = 0;
\end{equation}
\begin{equation}
\sum M_z = PR - M - Qsin30^\circ l = 0;
\end{equation}
\begin{equation}
\sum M_x = Y_Al - Psin60^\circ \cdot 0,5 \cdot 2l + 4Qsin30^\circ l - 5Y_Bl= 0;
\end{equation}
\begin{equation}
\sum M_y = -X_Al - Pcos60^\circ l + Qcos30^\circ cos 60^\circ l +5X_Bl = 0;
\end{equation}

Из уравнения $(25)$ получим:
\begin{equation}
Q = \frac{PR-M}{lsin30^\circ} = \frac{Pl}{2\cdot0,5l} = P.
\end{equation}

Подставим значение $Q$ в уравнение $(24)$:
\begin{equation}
Z_A = Pcos30^\circ cos60^\circ = 0,43P.
\end{equation}

Из уравнения $(22)$ выразим $X_A$ и подставим в уравнение $(27)$:
\begin{equation}
X_A = Pcos60^\circ - Qcos30^\circ sin60^\circ - X_B = -0,25P - X_B;
\end{equation}
\begin{equation}
-0,25Pl - X_Bl -0,5Pl +0,43Pl +5X_Bl = 0;
\end{equation}
\begin{equation}
X_B = 0,08P.
\end{equation}

Тогда
\begin{equation}
X_A = -0,33P.
\end{equation}

Из уравнения $(23)$ выразим $Y_A$ и подставим в уравнение $(26)$:
\begin{equation}
Y_A = Psin60^\circ + Qsin30^\circ - Y_B = 1,37P - Y_B;
\end{equation}
\begin{equation}
1,37Pl - Y_Bl - 0,87P +4\cdot0,5Pl -5Y_Bl = 0;
\end{equation}
\begin{equation}
Y_B = 0,42P.
\end{equation}

Тогда
\begin{equation}
Y_A = 0,95P.
\end{equation}

Тогда реакции подпятника $A$ и подшипника $B$ равны:
\begin{equation}
R_A = \sqrt{X^2_A + Y^2_A + Z^2_A} = \sqrt{(-0,33P)^2 + (0,95P)^2 + (0,43P)^2} = 1,09P;
\end{equation}
\begin{equation}
R_{B} = \sqrt{X^2_B + Y^2_B} = \sqrt{(0,08P)^2 + (0,42P)^2} = 0,43P.
\end{equation}

Так как значение составляющей реакции подпятника $A$ вдоль оси $Ox$ $X_A$ получилась в ходе расчетов отрицательной, то она должна быть направлена в противоположную сторону.

Вывод: в результате расчетов были получены значения сил реакции подпятника $A$ и подшипника $B$ и необходимой для равновесия силы $Q$, которые представлены в таблице 2.

\newpage

\begin{flushright}
  \textit{Таблица 2} \hspace*{0.4cm}
\end{flushright}

\vspace*{-0.9cm}

\begin{center}
  Значения сил реакции подпятника и подшипника и необходимой для равновесия силы
\end{center}

\begin{tabular}{|c|c|c|c|c|}
  \hline
  \rule{0pt}{1cm} № & \hspace*{0.5cm} Название величины \hspace*{0.5cm} & \makecell{Обозначение\\ величины} & Значение & \makecell{Единицы\\ измерения} \\[0.5cm]
  \hline
  \rule{0pt}{1cm} 1  & \makecell{Сила реакции \\ подпятника $A$} & $R_A$ & $1,09P$ & Н \\[0.5cm]
  \hline
  \rule{0pt}{1cm} 2  & \makecell{Сила реакции \\ подшипника $B$} & $R_B$ & $0,43P$ & Н \\[0.5cm]
  \hline
  \rule{0pt}{1cm} 2  & \makecell{Сила, необходимая\\для равновесия} & $Q$ & $P$ & Н \\[0.5cm]
  \hline
\end{tabular}



\end{document}

